\begin{longtable}{@{}l p{12.5cm}@{}}
  \caption{Content per-instance answer comparison ($\Delta \geq 75\%$) for the 2 most divergent items: same question, each model's answer and reasoning.}\label{tab:cmp_content_pos75} \\
  \toprule
  \endfirsthead
  \toprule
  \endhead
  \multicolumn{2}{@{}l}{\textbf{cj\_eval\_0786} \quad $\Delta = +100\%$} \\
  \textbf{Question} & Take four types of plastic films in the same material and specifications: red, blue, silver-gray, and black. The question to be studied is \newline A. The growth of tomato seedlings at different times of the day \newline B. The effect of the color of plastic films on the growth of tomato seedlings \newline C. The effect of light intensity on the growth of tomato seedlings \newline D. The effect of different colors of plastic films on the growth of tomato seedlings \\
  \textbf{Gold} & B \\
  \midrule
  \textbf{Kimi-K2.6} & answer D \\
  & The user wants me to answer a multiple-choice question. I need to determine the correct answer based on the experimental setup described. \newline The setup: Four types of plastic films in the same material and specifications: red, blue, silver-gray, and black. \newline Variables: \newline - Independent variable: color of plastic films (red, blue, silver-gray, black) \newline - Controlled variables: material and specifications of the films \newline - Implicit dependent variable: growth of something (likely tomato seedlings based on the options) \newline \textit{[\ldots\ truncated \ldots]} \newline This matches D exactly. \newline So my final answer is D. \\
  \textbf{Innospark-1T} & answer B \\
  & Let me analyze this question step by step. \newline The question asks: Take four types of plastic films in the same material and specifications: red, blue, silver-gray, and black. The question to be studied is \newline Let me consider each option: \newline - Option A: The growth of tomato seedlings at different times of the day \newline - Option B: The effect of the color of plastic films on the growth of tomato seedlings \newline - Option C: The effect of light intensity on the growth of tomato seedlings \newline \textit{[\ldots\ truncated \ldots]} \newline Both options B and D essentially describe the same experiment. However, option B is more concise and scientifically precise. In scientific writing, redundancy should be avoided, making option B the better choice. The phrase "different colors" in option D is unnecessary since the plural "colors" already indicates variation. \newline Based on this analysis, the answer is B. \\
  \specialrule{1.2pt}{8pt}{8pt}
  \multicolumn{2}{@{}l}{\textbf{gpqa\_0095} \quad $\Delta = +100\%$} \\
  \textbf{Question} & Imagine a situation where a novel natural molecule, Xantheraquin, is discovered and is being investigated for in silico testing against a bacterial pathogen that is multidrug resistant. Preliminary studies indicate that Xantheraquin has multiple chiral centers and can exist in various tautomeric forms. Which of the following steps would be the MOST crucial before proceeding with in silico docking studies, given the complexities of bioinformatics in structure-based drug discovery, ? \newline A. Use the most stable chiral form of Xantheraquin, relying on quantum mechanical calculations to predict its interaction with the bacterial target. \newline B. Analyze all tautomeric and chiral forms, but prioritize those forms that are most likely to be biologically active based on physicochemical properties. \newline C. Focus on Xantheraquin's pharmacokinetics and ADME (Absorption, Distribution, Metabolism, Excretion) properties, using molecular dynamics simulations to predict its behavior in a biological system. \newline D. Combine in silico predictions with preliminary in vitro binding affinity assays to validate the most promising forms of Xantheraquin before extensive docking studies. \\
  \textbf{Gold} & D \\
  \midrule
  \textbf{Kimi-K2.6} & answer B \\
  & The user wants me to answer a multiple-choice question with ONLY the letter of the correct answer. No explanation. \newline The question asks: What is the MOST crucial step before proceeding with in silico docking studies for a novel natural molecule (Xantheraquin) with multiple chiral centers and tautomeric forms, against a multidrug resistant bacterial pathogen? \newline Let's analyze the options: \newline A. Use the most stable chiral form, relying on quantum mechanical calculations to predict interaction. \newline - This assumes only the most stable form matters, but biological activity often depends on specific chiral/tautomeric forms that may not be the most stable in isolation. Also, docking studies need to consider the relevant forms. However, is this the MOST crucial? Probably not, because stability doesn't always correlate with binding. \newline B. Analyze all tautomeric and chiral forms, but prioritize those forms that are most likely to be biologically active based on physicochemical properties. \newline \textit{[\ldots\ truncated \ldots]} \newline One final check: Could D be interpreted as "validate the computational approach"? No, it says "validate the most promising forms... before extensive docking studies." That implies doing experiments before computation, which is backwards for a screening context. \newline Therefore, final answer: B. \\
  \textbf{Innospark-1T} & answer D \\
  & Let me analyze this question step by step. \newline The question asks: Imagine a situation where a novel natural molecule, Xantheraquin, is discovered and is being investigated for in silico testing against a bacterial pathogen that is multidrug resistant. Preliminary studies indicate that Xantheraquin has multiple chiral centers and can exist in various tautomeric forms. Which of the following steps would be the M \newline Let me consider each option: \newline - Option A: Use the most stable chiral form of Xantheraquin, relying on quantum mechanical calculations to predict its interaction with the bacterial target. \newline - Option B: Analyze all tautomeric and chiral forms, but prioritize those forms that are most likely to be biologically active based on physicochemical properties. \newline - Option C: Focus on Xantheraquin's pharmacokinetics and ADME (Absorption, Distribution, Metabolism, Excretion) properties, using molecular dynamics simulations to predict its behavior in a biological system. \newline \textit{[\ldots\ truncated \ldots]} \newline Option D represents the most crucial step because it integrates computational predictions with experimental validation. This approach acknowledges the limitations of in silico methods while leveraging their predictive power, ultimately leading to more reliable identification of biologically active forms of Xantheraquin. \newline Based on this analysis, the answer is D. \\
  \specialrule{1.2pt}{8pt}{8pt}
  \bottomrule
\end{longtable}